\chapter{Force Generation, Motion Control, and Conclusions}

\section{Chapter Objective and Interface from Flux Control}

This chapter develops the force-generation and motion-control stages using validated outputs from the flux-control stage. The objective is to evaluate how upstream identification and control decisions influence downstream force quality, tracking behavior, and operational outcomes in a traceable manner.

\section{Force-Generation Interface and Allocation Logic}

This section defines the force-generation interface from desired task-level force behavior to actuator-level command variables. The emphasis is on interface consistency, allocation logic, and implementation assumptions required for reproducible operation.

\section{Force- and Motion-Quality Criteria and Scheduling Context}

This section defines force- and motion-quality criteria used in this dissertation, including main-axis quality, cross-axis interaction behavior, tracking-level behavior, and scheduling-related response characteristics. The criteria are selected to support control-oriented interpretation and chapter-to-chapter comparability.

\section{Simulation and Experimental Validation of Force and Motion Outputs}

This section presents simulation and experimental validation in the same technical thread to assess force-generation and motion-control behavior under representative operating conditions. The goal is to verify whether force and motion output quality are consistent with the intended design chain from identification and flux control.

\section{Tracing Design Decisions from Identification to Flux, Force, and Motion Control}

This section closes the evidence chain by tracing how upstream modeling and control decisions propagate to force-level and motion-level outcomes through validated interfaces. The discussion focuses on evidence-supported links under tested conditions, and on how this traceability supports reproducible development.

\section{Conclusions and Future Directions}

This section summarizes the dissertation-level contribution as an automation-oriented and control-grounded development framework for Hall-sensor-based hexapole magnetic tweezers. It also outlines practical next steps for extending the framework in future studies while preserving traceability and reproducibility.
